\section{Pregunta N$^{\circ}$12\qquad Alejandro Escobar Mejia}

\begin{frame}
	\begin{enumerate}\setcounter{enumi}{11}
		\item

		      Considere la solución numérica del sistema de ecuaciones
		      lineales $Ax=b$ con
		      \begin{math}
			      A=
			      \begin{pmatrix}
				      1 & 2 \\
				      3 & 4
			      \end{pmatrix}.
		      \end{math}

		      \begin{enumerate}[a)]
			      \item

			            Encuentre el número de condición de $A$ en la norma
			            infinito.

			      \item

			            Encuentre el lado derecho $b$ y la perturbación
			            $\Delta b$ para el cual este número de condición se
			            alcanza realmente.

			      \item

			            Para el mismo $b$, encuentre una perturbación
			            $\Delta A$ para el cual este número de condición se
			            alcanza aproximadamente.
		      \end{enumerate}
	\end{enumerate}
	\begin{solution}
		\begin{enumerate}[a)]
			\item
			      \begin{math}
				      {\left\|A\right\|}_{\infty}=
				      \max\limits_{1\leq i\leq 2}
				      \left\{
				      \sum\limits_{j=1}^{2}\left|a_{ij}\right|
				      \right\}=
				      \max\left\{\left(\left|1\right|+\left|2\right|,\left|3\right|+\left|4\right|\right)\right\}=
				      7
			      \end{math}.
			      Si
			      \begin{math}
				      A^{-1}=
				      \begin{pmatrix}
					      -2  & 1    \\
					      1.5 & -0.5
				      \end{pmatrix}
			      \end{math},
			      entonces

			      \begin{math}
				      {\left\|A^{-1}\right\|}_{\infty}=
				      \max\left\{\left(\left|-2\right|+\left|1\right|,\left|1.5\right|+\left|-0.5\right|\right)\right\}=
				      3
			      \end{math}.
			      Por lo tanto,
			      \begin{math}
				      \operatorname{Cond}\left(A\right)=
				      7\times3=
				      21.
			      \end{math}

			\item

			      Sean
			      \begin{math}
				      b
				      =
				      \begin{pmatrix}
					      1 \\
					      -1
				      \end{pmatrix}
			      \end{math},
			      \begin{math}
				      Ax=
				      \begin{pmatrix}
					      1 & 2 \\
					      3 & 4
				      \end{pmatrix}
				      \begin{pmatrix}
					      x_{1} \\
					      x_{2}
				      \end{pmatrix}
				      =
				      \begin{pmatrix}
					      1 \\
					      -1
				      \end{pmatrix}
				      \begin{pmatrix}
					      x_{1} \\
					      x_{2}
				      \end{pmatrix}
				      =
				      \begin{pmatrix}
					      -3 \\
					      2
				      \end{pmatrix}
			      \end{math}.

			      \begin{math}
				      b+\delta b =
				      \begin{pmatrix}
					      1.1 \\
					      -1
				      \end{pmatrix}
			      \end{math}.

			      \begin{math}
				      x+\delta x =
				      \begin{pmatrix}
					      -3.2 \\
					      2.15
				      \end{pmatrix}
			      \end{math}

			      \begin{math}
				      \frac{{\left\|\Delta b\right\|}_{\infty}}{{\left\|b\right\|}_{\infty}}=
				      0.1
				      \frac{{\left\|\Delta x\right\|}_{\infty}}{{\left\|x\right\|}_{\infty}}=
				      0.06
			      \end{math}.

			      \begin{math}
				      \operatorname{Cond}\left(A\right)>\frac{0.06}{0.1}=0.6
			      \end{math}
		\end{enumerate}
	\end{solution}
\end{frame}